% A sample book made with the memoir class.
% Compile with: latexmk -pdf -outdir=build book.tex
\documentclass[a4paper,11pt,openright]{memoir}

\usepackage[T1]{fontenc}
\usepackage{lmodern}
\usepackage{amsmath}
\usepackage{graphicx} % the veelo chapter style needs \resizebox
\usepackage{tikz}     % for the gorilla picture in gorilla.tex
\usepackage{lipsum} % filler text for the sample
\usepackage[hidelinks]{hyperref}

% memoir numbers only chapters and sections by default.
% These two lines also number subsections, and show the levels down to
% subsubsection in the contents (the protocols are on the subsubsection level).
\setsecnumdepth{subsection}
\settocdepth{subsubsection}

% --- \topic: a heading that looks like \subsection, with its own counter ---
%
% Commands to change the number, used before \topic:
%   \skiptopic          jump one number (\skiptopic[2] jumps two numbers)
%   \topicnumber{10}    the next topic gets the number 10, then 11, 12
%   \topicnumber*{10b}  the next topic gets the text "10b", the count does not change
\newcounter{topic}                                    % one counter for the full book
\newcounter{topicid}                                  % hidden counter, never changed by hand
\makeatletter
\let\topic@pending\@empty                             % free text for the next topic
\let\topic@override\@empty                            % free text for the current topic
% number format: 1, 2, 3, or the free text
\renewcommand{\thetopic}{\ifx\topic@override\@empty\arabic{topic}\else\topic@override\fi}
\newcommand{\skiptopic}[1][1]{\addtocounter{topic}{#1}}
\newcommand{\topicnumber}{\@ifstar\topic@settext\topic@setnumber}
\newcommand{\topic@setnumber}[1]{\setcounter{topic}{\numexpr(#1)-1\relax}}
\newcommand{\topic@settext}[1]{\gdef\topic@pending{#1}}
% \topic* has no number, so it does not use the free text
\newcommand{\topic}{%
  \stepcounter{topicid}%
  \@ifstar{\topic@heading*}{\topic@usetext\topic@heading}}
\newcommand{\topic@usetext}{%
  \global\let\topic@override\topic@pending
  \global\let\topic@pending\@empty
  % a numbered topic with free text must not use a number from the count
  \ifx\topic@override\@empty\else
    \ifnum\c@secnumdepth>1 \addtocounter{topic}{-1}\fi
  \fi}
\newcommand{\topic@heading}{%
  \subsechook\@startsection{topic}{2}{\subsecindent}%
  {\beforesubsecskip}{\aftersubsecskip}{\normalfont\subsecheadstyle}}
\newcommand{\topicmark}[1]{}                          % no entry in the page header
\newcommand{\topicname}{Experiment}                   % the word before the number
% In the text: put \topicname before the number of a topic only.
\renewcommand{\@seccntformat}[1]{%
  \@ifundefined{#1@cntformat}{\csname the#1\endcsname\quad}{\csname #1@cntformat\endcsname}}
\newcommand{\topic@cntformat}{\topicname~\thetopic\quad}
% In the contents: use the subsection format, with the word before the number.
% The number box has space for the word, a number of two digits, and a small gap.
\newcommand{\l@topic}[2]{{%
  \settowidth{\cftsubsectionnumwidth}{\cftsubsectionfont\topicname~00\hspace{0.6em}}%
  \renewcommand{\cftsubsectionpresnum}{\topicname~}%
  % if the number is wider than the box (free text), make the box wider
  \renewcommand{\numberline}[1]{%
    \settowidth{\@tempdimb}{\@cftbsnum ##1\@cftasnum\hspace{0.6em}}%
    \ifdim\@tempdimb<\@tempdima \@tempdimb=\@tempdima \fi
    \hb@xt@\@tempdimb{\@cftbsnum ##1\@cftasnum\hfil}\@cftasnumb}%
  \l@subsection{#1}{#2}}}
\newcommand{\toclevel@topic}{2}                       % same level in PDF bookmarks
\renewcommand{\theHtopic}{\arabic{topicid}}           % unique link targets, also with equal numbers
\makeatother

% --- \protocol and \protocolsummary: headings on the subsubsection level ---
%
% The two commands use one counter. The title is optional:
%   \protocol                    gives "PROTOCOL 1"
%   \protocolsummary[The Title]  gives "PROTOCOL SUMMARY 2 The Title"
% Commands to change the number, used before one of the two commands:
%   \skipprotocol          jump one number (\skipprotocol[2] jumps two numbers)
%   \protocolnumber{10}    the next heading gets the number 10, then 11, 12
%   \protocolnumber*{10a}  the next heading gets the text "10a", the count does not change
\newcounter{protocol}                                 % one counter for the two commands
\newcounter{protocolid}                               % hidden counter, never changed by hand
\makeatletter
\newcommand{\protocolname}{PROTOCOL}
\newcommand{\protocolsummaryname}{PROTOCOL SUMMARY}
\let\protocol@pending\@empty                          % free text for the next heading
\let\protocol@override\@empty                         % free text for the current heading
\renewcommand{\theprotocol}{\ifx\protocol@override\@empty\arabic{protocol}\else\protocol@override\fi}
\renewcommand{\theHprotocol}{\arabic{protocolid}}     % unique link targets
\newcommand{\skipprotocol}[1][1]{\addtocounter{protocol}{#1}}
\newcommand{\protocolnumber}{\@ifstar\protocol@settext\protocol@setnumber}
\newcommand{\protocol@setnumber}[1]{\setcounter{protocol}{\numexpr(#1)-1\relax}}
\newcommand{\protocol@settext}[1]{\gdef\protocol@pending{#1}}
\newcommand{\protocol}[1][]{\protocol@heading{\protocolname}{#1}}
\newcommand{\protocolsummary}[1][]{\protocol@heading{\protocolsummaryname}{#1}}
% The word, the number, and the title are one text with normal spaces.
% There is no number box, thus there is no extra space after the number.
\newcommand{\protocol@text}[2]{#1~\theprotocol\if\relax\detokenize{#2}\relax\else\space #2\fi}
% The heading is a subsubsection without a number of its own (star form).
\newcommand{\protocol@heading}[2]{%
  \subsubsechook
  \@startsection{subsubsection}{3}{\subsubsecindent}%
    {\beforesubsubsecskip}{\aftersubsubsecskip}{\normalfont\subsubsecheadstyle}*%
    {\protocol@setup{#1}{#2}\protocol@text{#1}{#2}}}
% This code runs in the line of the heading: it sets the number, makes the
% target for \label and for links, and writes the line for the contents.
\newcommand{\protocol@setup}[2]{%
  \stepcounter{protocolid}%
  \global\let\protocol@override\protocol@pending
  \global\let\protocol@pending\@empty
  % a heading with free text must not use a number from the count
  \ifx\protocol@override\@empty\else\addtocounter{protocol}{-1}\fi
  \refstepcounter{protocol}%
  % the heading is in a group, thus make the data for \label global
  \global\let\@currentlabel\@currentlabel
  \global\let\@currentHref\@currentHref
  \addcontentsline{toc}{subsubsection}{\protocol@text{#1}{#2}}}
\makeatother

\chapterstyle{veelo}
\pagestyle{ruled}

\title{A Sample Book}
\author{Author Name}
\date{\today}

\begin{document}

\frontmatter
\maketitle
\clearpage
\tableofcontents

\mainmatter

\part{Foundations}

\chapter{First Steps}
\label{ch:first-steps}
This chapter shows the basic structure of a book. \lipsum[1]

\section{A Section}
\label{sec:a-section}
A section divides a chapter. \lipsum[2]

\subsection{A Subsection}
A subsection divides a section. \lipsum[3]

\subsection{A Second Subsection}
Text can contain a formula:
\begin{equation}
  \label{eq:euler}
  e^{i\pi} + 1 = 0
\end{equation}
Equation~\eqref{eq:euler} is in Section~\ref{sec:a-section}.

\topic{A First Topic}
\label{top:first}
A topic looks like a subsection, but it has its own counter. \lipsum[3]

\topic{A Second Topic}
Topic~\ref{top:first} is the topic before this one.

\section{A Second Section}
\lipsum[4]

\chapter{More Structure}
\label{ch:more-structure}
Chapter~\ref{ch:first-steps} starts on page~\pageref{ch:first-steps}. \lipsum[5]

\section{Lists}
\begin{itemize}
  \item A part contains chapters.
  \item A chapter contains sections.
  \item A section contains subsections.
\end{itemize}

\subsection{A Table}
Table~\ref{tab:levels} shows the levels of the book.

\begin{table}[ht]
  \centering
  \begin{tabular}{lll}
    \toprule
    Level & Command & Numbered \\
    \midrule
    Part       & \verb|\part|       & I, II \\
    Chapter    & \verb|\chapter|    & 1, 2 \\
    Section    & \verb|\section|    & 1.1 \\
    Subsection & \verb|\subsection| & 1.1.1 \\
    \bottomrule
  \end{tabular}
  \caption{The levels of the book structure.}
  \label{tab:levels}
\end{table}

\subsection{A Picture}
Figure~\ref{fig:gorilla} shows a gorilla. The picture is made with TikZ shapes.

\begin{figure}[ht]
  \centering
  \newcommand{\gorillascale}{0.8}
  % A gorilla made with TikZ shapes. Use: % A gorilla made with TikZ shapes. Use: % A gorilla made with TikZ shapes. Use: \input{gorilla} (needs \usepackage{tikz}).
% Change the size with \newcommand{\gorillascale}{0.8} before \input.
\providecommand{\gorillascale}{1}
\begin{tikzpicture}[scale=\gorillascale, line cap=round, line join=round]
  \colorlet{fur}{black!88}
  \colorlet{furlight}{black!72}
  \colorlet{skin}{black!52}
  \colorlet{skinlight}{black!40}

  % ground shadow
  \fill[black!12] (0,-3.25) ellipse [x radius=4.3, y radius=0.35];

  % shoulders and body
  \fill[fur] (0,0.7) ellipse [x radius=2.9, y radius=1.5];
  \fill[fur] (0,-0.9) ellipse [x radius=2.4, y radius=2.2];

  % chest
  \fill[furlight] (0,-0.1) ellipse [x radius=1.7, y radius=1.15];
  \draw[fur, line width=1.2pt] (0,0.95) -- (0,-1.1);
  \draw[fur, line width=1.2pt] (-1.2,-0.75) .. controls (-0.6,-1.15) and (0.6,-1.15) .. (1.2,-0.75);

  % feet
  \foreach \s in {-1,1} {
    \begin{scope}[xscale=\s]
      \fill[fur] (1.25,-2.55) ellipse [x radius=1.15, y radius=0.7];
      \fill[skin] (1.35,-2.85) ellipse [x radius=0.85, y radius=0.33];
      \foreach \t in {0.75,1.1,1.45,1.8}
        \draw[fur, line width=0.8pt] (\t,-2.72) -- (\t,-3.05);
    \end{scope}
  }

  % arms and fists
  \foreach \s in {-1,1} {
    \begin{scope}[xscale=\s]
      \fill[fur] (1.7,1.9)
        .. controls (3.7,1.9) and (4.3,-0.8) .. (3.9,-2.8)
        -- (2.5,-2.8)
        .. controls (2.7,-1.5) and (2.5,-0.4) .. (1.9,0.2) -- cycle;
      \draw[furlight, line width=1.2pt] (2.75,0.9) .. controls (3.2,0.0) and (3.3,-1.0) .. (3.2,-1.9);
      \fill[skin] (3.2,-2.9) ellipse [x radius=0.85, y radius=0.36];
      \foreach \t in {2.75,3.05,3.35,3.65}
        \draw[fur, line width=0.8pt] (\t,-2.72) -- (\t,-3.1);
    \end{scope}
  }

  % head: ears, skull, crest
  \foreach \s in {-1,1} {
    \fill[fur] (\s*1.32,2.35) circle [radius=0.36];
    \fill[skin] (\s*1.36,2.35) circle [radius=0.19];
  }
  \fill[fur] (0,2.3) ellipse [x radius=1.35, y radius=1.35];
  \fill[fur] (-1.0,3.0) .. controls (-0.6,4.1) and (0.6,4.1) .. (1.0,3.0) -- cycle;

  % face
  \fill[skin] (0,2.05) ellipse [x radius=0.98, y radius=0.92];
  \fill[skinlight] (0,1.5) ellipse [x radius=0.8, y radius=0.55];

  % brow
  \fill[fur] (-1.05,2.5) .. controls (-0.5,2.95) and (0.5,2.95) .. (1.05,2.5)
    .. controls (0.5,2.45) and (0.2,2.3) .. (0,2.3)
    .. controls (-0.2,2.3) and (-0.5,2.45) .. cycle;

  % eyes
  \foreach \s in {-1,1} {
    \fill[fur] (\s*0.42,2.2) ellipse [x radius=0.24, y radius=0.17];
    \fill[brown!70!orange] (\s*0.42,2.18) circle [radius=0.11];
    \fill[black] (\s*0.42,2.18) circle [radius=0.055];
    \fill[white] (\s*0.42+0.03,2.21) circle [radius=0.02];
  }

  % nose and mouth
  \fill[skin] (0,1.72) ellipse [x radius=0.42, y radius=0.24];
  \foreach \s in {-1,1}
    \fill[black] (\s*0.17,1.7) ellipse [x radius=0.12, y radius=0.08];
  \draw[fur, line width=1pt] (-0.5,1.36) .. controls (-0.2,1.24) and (0.2,1.24) .. (0.5,1.36);
\end{tikzpicture}
 (needs \usepackage{tikz}).
% Change the size with \newcommand{\gorillascale}{0.8} before \input.
\providecommand{\gorillascale}{1}
\begin{tikzpicture}[scale=\gorillascale, line cap=round, line join=round]
  \colorlet{fur}{black!88}
  \colorlet{furlight}{black!72}
  \colorlet{skin}{black!52}
  \colorlet{skinlight}{black!40}

  % ground shadow
  \fill[black!12] (0,-3.25) ellipse [x radius=4.3, y radius=0.35];

  % shoulders and body
  \fill[fur] (0,0.7) ellipse [x radius=2.9, y radius=1.5];
  \fill[fur] (0,-0.9) ellipse [x radius=2.4, y radius=2.2];

  % chest
  \fill[furlight] (0,-0.1) ellipse [x radius=1.7, y radius=1.15];
  \draw[fur, line width=1.2pt] (0,0.95) -- (0,-1.1);
  \draw[fur, line width=1.2pt] (-1.2,-0.75) .. controls (-0.6,-1.15) and (0.6,-1.15) .. (1.2,-0.75);

  % feet
  \foreach \s in {-1,1} {
    \begin{scope}[xscale=\s]
      \fill[fur] (1.25,-2.55) ellipse [x radius=1.15, y radius=0.7];
      \fill[skin] (1.35,-2.85) ellipse [x radius=0.85, y radius=0.33];
      \foreach \t in {0.75,1.1,1.45,1.8}
        \draw[fur, line width=0.8pt] (\t,-2.72) -- (\t,-3.05);
    \end{scope}
  }

  % arms and fists
  \foreach \s in {-1,1} {
    \begin{scope}[xscale=\s]
      \fill[fur] (1.7,1.9)
        .. controls (3.7,1.9) and (4.3,-0.8) .. (3.9,-2.8)
        -- (2.5,-2.8)
        .. controls (2.7,-1.5) and (2.5,-0.4) .. (1.9,0.2) -- cycle;
      \draw[furlight, line width=1.2pt] (2.75,0.9) .. controls (3.2,0.0) and (3.3,-1.0) .. (3.2,-1.9);
      \fill[skin] (3.2,-2.9) ellipse [x radius=0.85, y radius=0.36];
      \foreach \t in {2.75,3.05,3.35,3.65}
        \draw[fur, line width=0.8pt] (\t,-2.72) -- (\t,-3.1);
    \end{scope}
  }

  % head: ears, skull, crest
  \foreach \s in {-1,1} {
    \fill[fur] (\s*1.32,2.35) circle [radius=0.36];
    \fill[skin] (\s*1.36,2.35) circle [radius=0.19];
  }
  \fill[fur] (0,2.3) ellipse [x radius=1.35, y radius=1.35];
  \fill[fur] (-1.0,3.0) .. controls (-0.6,4.1) and (0.6,4.1) .. (1.0,3.0) -- cycle;

  % face
  \fill[skin] (0,2.05) ellipse [x radius=0.98, y radius=0.92];
  \fill[skinlight] (0,1.5) ellipse [x radius=0.8, y radius=0.55];

  % brow
  \fill[fur] (-1.05,2.5) .. controls (-0.5,2.95) and (0.5,2.95) .. (1.05,2.5)
    .. controls (0.5,2.45) and (0.2,2.3) .. (0,2.3)
    .. controls (-0.2,2.3) and (-0.5,2.45) .. cycle;

  % eyes
  \foreach \s in {-1,1} {
    \fill[fur] (\s*0.42,2.2) ellipse [x radius=0.24, y radius=0.17];
    \fill[brown!70!orange] (\s*0.42,2.18) circle [radius=0.11];
    \fill[black] (\s*0.42,2.18) circle [radius=0.055];
    \fill[white] (\s*0.42+0.03,2.21) circle [radius=0.02];
  }

  % nose and mouth
  \fill[skin] (0,1.72) ellipse [x radius=0.42, y radius=0.24];
  \foreach \s in {-1,1}
    \fill[black] (\s*0.17,1.7) ellipse [x radius=0.12, y radius=0.08];
  \draw[fur, line width=1pt] (-0.5,1.36) .. controls (-0.2,1.24) and (0.2,1.24) .. (0.5,1.36);
\end{tikzpicture}
 (needs \usepackage{tikz}).
% Change the size with \newcommand{\gorillascale}{0.8} before \input.
\providecommand{\gorillascale}{1}
\begin{tikzpicture}[scale=\gorillascale, line cap=round, line join=round]
  \colorlet{fur}{black!88}
  \colorlet{furlight}{black!72}
  \colorlet{skin}{black!52}
  \colorlet{skinlight}{black!40}

  % ground shadow
  \fill[black!12] (0,-3.25) ellipse [x radius=4.3, y radius=0.35];

  % shoulders and body
  \fill[fur] (0,0.7) ellipse [x radius=2.9, y radius=1.5];
  \fill[fur] (0,-0.9) ellipse [x radius=2.4, y radius=2.2];

  % chest
  \fill[furlight] (0,-0.1) ellipse [x radius=1.7, y radius=1.15];
  \draw[fur, line width=1.2pt] (0,0.95) -- (0,-1.1);
  \draw[fur, line width=1.2pt] (-1.2,-0.75) .. controls (-0.6,-1.15) and (0.6,-1.15) .. (1.2,-0.75);

  % feet
  \foreach \s in {-1,1} {
    \begin{scope}[xscale=\s]
      \fill[fur] (1.25,-2.55) ellipse [x radius=1.15, y radius=0.7];
      \fill[skin] (1.35,-2.85) ellipse [x radius=0.85, y radius=0.33];
      \foreach \t in {0.75,1.1,1.45,1.8}
        \draw[fur, line width=0.8pt] (\t,-2.72) -- (\t,-3.05);
    \end{scope}
  }

  % arms and fists
  \foreach \s in {-1,1} {
    \begin{scope}[xscale=\s]
      \fill[fur] (1.7,1.9)
        .. controls (3.7,1.9) and (4.3,-0.8) .. (3.9,-2.8)
        -- (2.5,-2.8)
        .. controls (2.7,-1.5) and (2.5,-0.4) .. (1.9,0.2) -- cycle;
      \draw[furlight, line width=1.2pt] (2.75,0.9) .. controls (3.2,0.0) and (3.3,-1.0) .. (3.2,-1.9);
      \fill[skin] (3.2,-2.9) ellipse [x radius=0.85, y radius=0.36];
      \foreach \t in {2.75,3.05,3.35,3.65}
        \draw[fur, line width=0.8pt] (\t,-2.72) -- (\t,-3.1);
    \end{scope}
  }

  % head: ears, skull, crest
  \foreach \s in {-1,1} {
    \fill[fur] (\s*1.32,2.35) circle [radius=0.36];
    \fill[skin] (\s*1.36,2.35) circle [radius=0.19];
  }
  \fill[fur] (0,2.3) ellipse [x radius=1.35, y radius=1.35];
  \fill[fur] (-1.0,3.0) .. controls (-0.6,4.1) and (0.6,4.1) .. (1.0,3.0) -- cycle;

  % face
  \fill[skin] (0,2.05) ellipse [x radius=0.98, y radius=0.92];
  \fill[skinlight] (0,1.5) ellipse [x radius=0.8, y radius=0.55];

  % brow
  \fill[fur] (-1.05,2.5) .. controls (-0.5,2.95) and (0.5,2.95) .. (1.05,2.5)
    .. controls (0.5,2.45) and (0.2,2.3) .. (0,2.3)
    .. controls (-0.2,2.3) and (-0.5,2.45) .. cycle;

  % eyes
  \foreach \s in {-1,1} {
    \fill[fur] (\s*0.42,2.2) ellipse [x radius=0.24, y radius=0.17];
    \fill[brown!70!orange] (\s*0.42,2.18) circle [radius=0.11];
    \fill[black] (\s*0.42,2.18) circle [radius=0.055];
    \fill[white] (\s*0.42+0.03,2.21) circle [radius=0.02];
  }

  % nose and mouth
  \fill[skin] (0,1.72) ellipse [x radius=0.42, y radius=0.24];
  \foreach \s in {-1,1}
    \fill[black] (\s*0.17,1.7) ellipse [x radius=0.12, y radius=0.08];
  \draw[fur, line width=1pt] (-0.5,1.36) .. controls (-0.2,1.24) and (0.2,1.24) .. (0.5,1.36);
\end{tikzpicture}

  \caption{A gorilla made with TikZ.}
  \label{fig:gorilla}
\end{figure}

\skiptopic
\topic{A Topic in Chapter Two}
The counter of the topics does not start again in a new chapter.
The command \verb|\skiptopic| before this topic jumps the number 3.

\part{Practice}

\chapter{A Longer Example}
\lipsum[6-7]

\section{Method}
\lipsum[8]

\subsection{Setup}
\lipsum[9]

\subsection{Results}
\lipsum[10]

\topic{A Topic in Part Two}
This topic is in the second part of the book. Its number continues from the first part.

\protocol
\label{pro:first}
A protocol without a title. \lipsum[2][1-2]

\protocolsummary
A protocol summary without a title. It uses the same counter as the protocols.

\protocol[Preparation of the Samples]
\label{pro:samples}
A protocol with a title. There is a normal space between the number and the title.

\protocolsummary[Results of the Preparation]
A protocol summary with a title.

\protocolnumber{10}
\protocol[A Protocol with a Set Number]
The command \verb|\protocolnumber{10}| sets the number of this protocol.

\protocolnumber*{10a}
\protocolsummary[A Summary with Free Text]
\label{pro:free}
The command \verb|\protocolnumber*{10a}| sets free text as the number.

\skipprotocol
\protocol
\label{pro:last}
The command \verb|\skipprotocol| jumps the number 11.
The labels give the numbers \ref{pro:first}, \ref{pro:samples}, \ref{pro:free}, and~\ref{pro:last}.

\section{Discussion}
\lipsum[11]

\chapter{Conclusion}
\lipsum[12]

\section{Summary}
\lipsum[13]

\topicnumber{10}
\topic{A Topic with a Set Number}
\label{top:set}
The command \verb|\topicnumber{10}| sets the number of this topic.

\topicnumber*{10b}
\topic{A Topic with Free Text}
\label{top:free}
The command \verb|\topicnumber*{10b}| sets free text as the number.
The count does not change.

\topic{A Topic after the Changes}
\label{top:after}
The count continues. The three topics have the numbers
\ref{top:set}, \ref{top:free}, and~\ref{top:after}.

\subsection{Next Steps}
Change the text, add a chapter, and compile the book again.

\backmatter

\chapter{Afterword}
The back matter has chapters without numbers. \lipsum[14]

\end{document}
